% Build: cd docs/reports/reconstruction/methods && latexmk -xelatex cross_experiment_ms_handling_20260427_zh.tex
\documentclass[a4paper, 11pt]{article}

\usepackage[margin=2.4cm]{geometry}
\usepackage{ctex}            % 中文支持
\usepackage{amsmath,amssymb}
\usepackage{booktabs}
\usepackage{multirow}
\usepackage{siunitx}
\usepackage{xcolor}
\usepackage{hyperref}
\usepackage{enumitem}
\usepackage{float}
\usepackage{graphicx}
\usepackage{longtable}
\usepackage{tabularx}

\definecolor{goodgreen}{RGB}{34,139,34}
\definecolor{warnred}{RGB}{200,50,50}

\hypersetup{
  colorlinks=true,
  linkcolor=blue!60!black,
  citecolor=blue!60!black,
  urlcolor=blue!60!black,
}

\title{\textbf{同能区带电粒子径迹重建中的 \\ 多重散射处理对照 \\ （SAMURAI 质子/碎片臂 + GSI）}}
\author{TBT \\ RIKEN 仁科中心 / DPOL 实验}
\date{2026 年 4 月 27 日}

\begin{document}
\maketitle
\tableofcontents
\clearpage

\section{摘要 TL;DR}
% [content from Task 10 — written LAST]

\section{为什么这件事重要}
% [content from Task 5]

\section{硬件 / 物质预算切面}
% [content from Task 2]

\subsection{磁谱仪腔：真空 vs 空气}
\subsection{出射窗}
\subsection{径迹探测器气体}
\subsection{束流管 / 上游靶室}
\subsection{跨实验硬件对照表}

\section{算法切面}
% [content from Task 3]

\subsection{直线 / RK 反推：无显式 MS 项}
\subsection{Kalman 滤波 + MS 过程噪声协方差}
\subsection{GEANE / RKF 带 MS 自适应步长}
\subsection{查表 / 反卷积修正（事后）}
\subsection{跨实验算法对照表}

\section{模拟 / 校准切面}
% [content from Task 4]

\subsection{所选 MC 框架}
\subsection{MS 模型调谐方法}
\subsection{残差反卷积 / 事后 MS 修正}
\subsection{验证观测量}
\subsection{跨实验模拟/校准对照表}

\section{逐实验事实卡}
% [content from Task 6]

\subsection{SAMURAI 质子臂 (PDC1/PDC2)}
\subsection{SAMURAI 碎片臂 (FDC1/FDC2)}
\subsection{R3B / LAND (GSI/FAIR)}
\subsection{ALADIN-EOS (GSI)}
\subsection{FOPI (GSI)}
\subsection{HADES (GSI)}

\section{对我们 \texorpdfstring{$\sigma_y$}{sigma\_y} 缺口的启示}
% [content from Task 7]

\subsection{硬件路径}
\subsection{算法路径}
\subsection{模拟/校准路径}
\subsection{组合路径}

\section{开放问题 / 待核证}
% [content from Task 8]

\section{参考文献}
% [content from Task 9]

\subsection{SAMURAI 质子臂 (PDC1/PDC2)}
\subsection{SAMURAI 碎片臂 (FDC1/FDC2)}
\subsection{R3B / LAND (GSI/FAIR)}
\subsection{ALADIN-EOS (GSI)}
\subsection{FOPI (GSI)}
\subsection{HADES (GSI)}
\subsection{交叉算法（GENFIT, GEANE, MS 模型）}
\subsection{交叉 MC 框架}

\end{document}
